% litreview_revised.tex — Full proposal draft, Week 8
% Revised from litreview.tex (Week 5). Minor edits for transition into Methods.

\section{Literature Review}

Understanding individual differences in attention has long been a goal of cognitive neuroscience, but studying it rigorously requires measuring the brain. Attention and working memory are among the most consequential cognitive capacities humans possess. They predict academic achievement, occupational performance, and vulnerability to clinical conditions including ADHD, with large and reliable effect sizes across the population \citep{kofler2013}. These capacities are also not randomly distributed. Sustained attention and executive function are meaningfully associated with socioeconomic status, with children from lower-income backgrounds showing consistent disadvantages in working memory and attentional control that are reflected not only in behavior but in the structure and function of prefrontal and hippocampal brain regions \citep{farah2006, noble2015}. The implication is that the populations most affected by attentional difficulties are also the most likely to be excluded from the scientific infrastructure that studies those difficulties. fMRI scanners, the dominant tool for imaging attention in the brain, are expensive, physically fixed, and concentrated in wealthy, Western, academic institutions, which means that research in low- and middle-income countries, non-laboratory community settings, and resource-limited clinical environments is structurally underrepresented \citep{burns2019}. The present project is organized around the question of whether that constraint can be relaxed through the combination of portable neuroimaging and naturalistic experimental design.

Functional near-infrared spectroscopy (fNIRS) is the most promising candidate for portable neuroimaging. It works by shining near-infrared light (typically 650 to 900 nm) through the scalp and detecting changes in the absorption of oxygenated (HbO) and deoxygenated (HbR) hemoglobin in cortical tissue, which is the same hemodynamic signal that fMRI measures, though restricted to cortical surface regions accessible to light \citep{kamran2016}. Unlike fMRI, fNIRS requires no magnetic field or radio-frequency pulses, tolerates participant movement, and can be deployed in any setting without special infrastructure. These properties have allowed it to reach populations and contexts previously inaccessible to neuroimaging. \citet{burns2019} demonstrated that portable fNIRS could measure meaningful individual differences in prefrontal responses to persuasive messages in a native Arabic-speaking sample recruited in Jordan, a population almost entirely absent from social neuroscience, and replicated established findings about medial and ventromedial prefrontal cortex involvement in persuasion. This work directly establishes that fNIRS prefrontal signals are sensitive enough to detect socially meaningful individual differences outside laboratory settings, which is a core assumption of the present project. The limitation fNIRS cannot overcome is depth. It is physically incapable of measuring medial prefrontal cortex, subcortical structures, or any brain region far from the scalp-mounted optodes. Because attention depends on a distributed network that includes the hippocampus, thalamus, basal ganglia, and portions of cingulate cortex, none of which fNIRS can reach, a portable surface-only recording cannot by itself provide the kind of whole-brain picture that motivates research on neural individual differences.

The key insight motivating this project is that the limitation of surface-only measurement does not have to be fatal. Deep brain activity can be inferred from cortical activity through learned functional connectivity relationships. \citet{liu2015} provided early proof of concept by scanning 17 participants simultaneously with fNIRS and fMRI during multiple cognitive tasks and showing that activity in deep structures, including the insula, hippocampus, and caudate, could be significantly predicted from surface cortical fNIRS signals, with top-quintile predictions achieving correlations of approximately $r = 0.70$. The neural basis for this inferential approach is well established: deep and subcortical regions are functionally coupled with surface cortical regions, so their activity tends to rise and fall together. A computational model that has learned what the fMRI signal looks like when the cortical signal looks a certain way can exploit this coupling. The limitation in Liu et al.'s work was the requirement that the same individuals be scanned simultaneously in both modalities, which is expensive, logistically demanding, and ultimately defeats the purpose of building a portable tool. The methodological breakthrough that resolved this came from naturalistic paradigms, which provide a way to align signals across separately scanned participant groups.

The logic of naturalistic alignment rests on a substantial literature on inter-subject correlation (ISC), which refers to the phenomenon whereby independently scanned individuals who experience the same naturalistic stimulus produce temporally synchronized brain responses across a wide network of regions \citep{hasson2004, hasson2010, nastase2019}. When people watch the same film, their auditory cortices respond to dialogue in near-lockstep, their prefrontal regions engage with narrative tension at similar moments, and their default mode networks activate and deactivate in patterns that are recognizably consistent across individuals \citep{hasson2010}. This shared temporal structure is largely absent during structured cognitive tasks, where individual differences in strategy, arousal, and performance introduce person-specific variance that overwhelms the shared signal \citep{nastase2019}. \citet{finn2021} showed that this property of naturalistic stimuli has direct practical consequences: functional connectivity measured during movie watching yields substantially more accurate predictions of trait-level behavioral phenotypes than connectivity measured during rest, with highly social clips producing reliable predictions in under three minutes of data. Naturalistic viewing, in other words, does not merely produce strong shared responses. It also amplifies the behaviorally relevant structure of individual differences in brain organization. For this project, the relevant implication is that a movie-watching session provides the ideal calibration paradigm because the shared temporal structure of the film creates a strong ISC signal that makes it possible to train a cross-dataset mapping between fNIRS and fMRI participants who were never scanned together, while the within-person variability around that shared signal may encode individual cognitive profiles.

\citet{gao2025} operationalized this insight in the aPCR framework. By training an adapted principal component regression model on fNIRS and fMRI data collected during naturalistic TV episode viewing, across separate participant groups who watched the same content, they demonstrated that whole-brain fMRI dynamics could be significantly predicted from prefrontal surface signals alone in 66 of 122 brain regions ($q < 0.05$), including subcortical and medial regions physically inaccessible to fNIRS. The model also generalized to an independent naturalistic dataset recorded during a different TV series, establishing that the learned cross-modal mapping captures stimulus-general structure rather than overfitting to one particular narrative \citep{gao2025}. The Gao et al.\ framework was designed and tested entirely within naturalistic paradigms, with both training and test data coming from movie or TV episode viewing. The present project takes the harder step of asking whether a model calibrated on naturalistic data can transfer to structured cognitive tasks, where the stimulus-driven ISC that supported model training is largely absent. This is the key untested assumption, and the answer to it determines whether naturalistic calibration can serve as a general-purpose tool for portable neurocognitive assessment.

Evaluating whether the cross-task transfer is meaningful requires a framework for what meaningful brain predictions during attention tasks should look like. Two intersecting literatures are relevant here. First, there is strong evidence that the frontoparietal network, and particularly the dorsal attention network (DAN) and the cognitive control network (CONT), functions as a set of flexible hubs whose global connectivity patterns shift more dramatically across task demands than any other network in the brain \citep{cole2013}. In a study spanning 64 different cognitive tasks, \citet{cole2013} found that frontoparietal regions could identify the current task from their connectivity patterns alone, and that this held even for tasks the person had never practiced before, which suggests that these hubs encode task demands in a generalizable format. Because the fNIRS cap used in this project covers precisely these prefrontal and lateral parietal regions, the model has direct access to the surface signals most likely to encode task-relevant structure. Building on this, \citet{braun2015} showed in 344 participants that how much frontal networks reorganized their communication patterns during the n-back task predicted both working memory accuracy and scores on external neuropsychological measures of cognitive flexibility, suggesting that individual differences in frontal system reconfiguration during a cognitive challenge carry behavioral consequences that generalize well beyond the task itself. This finding bears directly on Phase 5 of the present project: if the aPCR-predicted brain patterns preserve this frontoparietal reconfiguration structure, they should correlate with behavioral performance on the n-back and CPT. More broadly, \citet{tavor2016} demonstrated in the Human Connectome Project that task-evoked brain activity during seven different cognitive paradigms could be accurately predicted from resting-state connectivity alone, without any task data from the individuals being predicted. The parallel to the present project is close. Just as resting-state connectivity encodes something like a neural blueprint for how a person will respond to cognitive demands, a naturalistic fNIRS session may encode enough of an individual's brain organization to allow meaningful predictions of their task-evoked dynamics.

The second relevant literature addresses how behavioral performance during cognitive tasks can be used to index the quality of attention regulation. Reaction time variability (RTV) refers to the trial-to-trial inconsistency in how quickly a person responds across a sustained performance task, and it is among the most robust behavioral markers of attentional control. A meta-analysis of 319 studies by \citet{kofler2013} found that individuals with ADHD showed significantly greater RTV than typically developing peers (Hedges' $g = 0.76$ for children, $g = 0.46$ for adults). This elevation was more discriminating than mean response speed, persisted across task types including the n-back and CPT, and was substantially reduced by stimulant treatment ($g = -0.74$) but unaffected by behavioral interventions. The choice of RTV as the primary individual-differences index in this project, rather than task accuracy or mean response time, is deliberate. Accuracy is susceptible to floor and ceiling effects: in healthy young adults, 1-back performance is near-perfect and fails to differentiate individuals. Mean reaction time, for its part, conflates speed-accuracy tradeoffs with attentional stability. RTV, by contrast, captures moment-to-moment fluctuations in the efficiency of cognitive processing and remains discriminating even when average performance is high, because it reflects the regularity of attentional engagement across trials rather than its central tendency. Critically, RTV is not a marker only of clinical disorder. \citet{macdonald2009} reviewed evidence suggesting that intraindividual variability in cognitive performance tracks the health and efficiency of underlying neural systems, specifically the stability of dopaminergic neuromodulation, the integrity of white matter connectivity, and the regulatory strength of prefrontal cortex, and that these relationships hold across the full range of healthy adults. The theoretical account connecting RTV to brain dynamics was provided by \citet{sonugabarke2007}, who proposed the default mode interference hypothesis. During goal-directed tasks, the default mode network (DMN) is normally suppressed to allow task-relevant networks to operate without competition. When this suppression fails periodically, the DMN briefly reactivates and disrupts processing of the current stimulus, producing an attentional lapse that manifests as an unusually slow response on that trial. High RTV is, on this account, a behavioral fingerprint of poor DMN suppression. The DMN is also one of the networks where the aPCR model shows strong fNIRS-to-fMRI predictions (7 of 26 ROIs significant), which creates a testable prediction for the present project: participants with higher RTV should show more intrusion of rest-like, DMN-active dynamics during their task fNIRS recordings, making their task signals structurally different from movie-watching signals and potentially altering the model's ability to generate meaningful task-period predictions.

Taken together, these three bodies of literature converge on the design of the present project and jointly justify its specific choices. The fNIRS-to-fMRI translation literature establishes the computational feasibility of cross-modal inference from cortical surface signals. The naturalistic paradigm literature explains why a movie provides the ideal training context: it generates the inter-subject correlation necessary for cross-dataset model training and amplifies individual variation in brain organization in ways that resting-state recordings do not \citep{finn2021, gao2025, hasson2010}. The RTV literature links a simple behavioral measure, one that requires no imaging at all, to the neural dynamics the model is trying to predict, providing a concrete individual-differences target for evaluating whether the cross-task predictions are behaviorally meaningful \citep{kofler2013}. The equity framing is not rhetorical. The same reason the project uses movie-watching rather than fMRI, namely that it requires no scanner, is the reason the approach could matter beyond this study. If a brief, passive, naturalistic fNIRS session produces predictions that correlate with attentional performance, then the field has a candidate method for conducting whole-brain-equivalent attention research in schools, clinics, and communities where MRI access is not available. These are precisely the settings and populations that the existing literature, with its structural dependence on scanner infrastructure, has historically excluded \citep{burns2019, farah2006, noble2015}. What follows describes how the two datasets are combined and analyzed to test whether that candidate method delivers.
